


\section{Základné pojmy}
V tejto kapitole sa zameriame na všeobecné pojmy a technológie, 
ktoré sú kľúčové pre správne pochopenie problematiky vývoja mobilných aplikácií a implementácie odporúčacích systémov.

\iffalse
\subsection{Architektúra a komunikácia v distribuovaných systémoch}
Základom modernej sieťovej komunikácie je efektívne rozdelenie úloh medzi jednotlivé uzly v sieti. Tento prístup umožňuje lepšiu správu zdrojov a vysokú škálovateľnosť aplikácií.


\subsubsection{Klient-server architektúra}
Architektúra klient-server predstavuje systém, ktorý vykonáva funkcie klienta aj servera s cieľom podporiť zdieľanie informácií medzi nimi. Tento model umožňuje viacerým používateľom pristupovať k rovnakej databáze v rovnakom čase, pričom databáza slúži na ukladanie veľkého množstva informácií \cite{sulyman2014}.
(Zdroj: \url{https://www.researchgate.net/profile/Shakirat-Sulyman/publication/271295146_Client-Server_Model/links/5864e11308ae8fce490c1b01/Client-Server-Model.pdf})

\begin{itemize}
    \item \textbf{Server:} Predstavuje počítač alebo softvérový systém, ktorý poskytuje dáta, zdroje alebo služby iným počítačom v sieti, ktoré označujeme ako klienti. \\ (Zdroj: \url{https://en.wikipedia.org/wiki/Server_(computing)})
    \item \textbf{Klient:} Ide o kus počítačového hardvéru alebo softvéru, ktorý pristupuje k službe sprístupnenej serverom v rámci modelu klient-server v počítačových sieťach. \\ (Zdroj: \url{https://en.wikipedia.org/wiki/Client_(computing)})
\end{itemize}


\subsubsection{Komunikačné rozhrania a protokoly (REST API a HTTP)}
Na to, aby si klient a server v distribuovanom systéme rozumeli, potrebujú spoločný komunikačný štýl a protokol. V súčasnosti je najrozšírenejším štandardom rozhranie \foreignlanguage{english}{Representational State Transfer} (REST API). Ide o aplikačné programové rozhranie, ktoré vyhovuje princípom architektonického štýlu REST. Jeho podstata spočíva v práci so zdrojmi, ktorými môže byť akákoľvek informácia, napríklad používateľ, produkt alebo dokument \cite{google_cloud_rest}.

Celý proces komunikácie prebieha v dvoch základných krokoch:
\begin{enumerate}
    \item \textbf{Klient odošle požiadavku:} Mobilná aplikácia iniciuje požiadavku na vykonanie akcie smerom k API. Táto požiadavka obsahuje metódu HTTP, hlavičky s metadátami a prípadne telo s dátami.
    \item \textbf{Server spracuje požiadavku:} Server prijme a overí požiadavku (autentifikácia a autorizácia). Následne ju spracuje, čo môže zahŕňať načítanie dát z databázy alebo vytvorenie nového zdroja. \\ (Zdroj: \url{https://cloud.google.com/discover/what-is-rest-api})
\end{enumerate}

Samotným nosným médiom pre túto komunikáciu je protokol \foreignlanguage{english}{Hypertext Transfer Protocol} (HTTP). Ide o základný aplikačný protokol, ktorý umožňuje komunikáciu a prenos dát v prostredí celosvetovej siete. V architektúre klient-server funguje ako „jazyk“, ktorým sa tieto dve strany dorozumievajú. \\ (Zdroj: \url{https://gemini.google.com/app/6ed6e44697b48104} -- konzultácia 29.03.2026)

\subsubsection{Cloudové riešenia a služby}
Cloud môžeme prirovnať k virtuálnemu priestoru, prostredníctvom ktorého používatelia vytvárajú, zdieľajú a narábajú s obrovským množstvom informácií rôzneho typu. Ide o výpočtové zdroje dostupné na vyžiadanie (\foreignlanguage{english}{on-demand}), ku ktorým má používateľ prístup cez internet \cite{eset_cloud}.

Dáta sa ukladajú na cloudové servery umiestnené v dátových centrách po celom svete a sú plne riadené ich poskytovateľom. Za poplatok môže zákazník na týchto serveroch prevádzkovať svoj softvér, využívať ich výpočtový výkon či úložisko na ukladanie vlastných dát. \\ (Zdroj: \url{https://bezpecnenanete.eset.com/sk/it-bezpecnost/co-je-to-cloud-a-ako-funguje/})

\subsection{Správa a ukladanie dát}
Databáza funguje ako digitálne úložisko, ktoré organizuje dáta potrebné pre chod aplikácie. Podľa spôsobu ukladania a výkonu ich najčastejšie delíme na:

\subsubsection{Relačné databázy (SQL)}
Dáta sú v tomto modeli organizované do \textbf{tabuliek} s riadkami a stĺpcami. Medzi najpoužívanejšie systémy podľa portálu Webglobe \cite{webglobe_databaza} patria:
\begin{itemize}
    \item \textbf{MySQL:} Jednoduchá a rýchla databáza, hojne využívaná pri menších až stredne veľkých aplikáciách.
    \item \textbf{MariaDB:} Ponúka vyšší výkon a rozšírené funkcie, je vhodná aj pre náročnejšie aplikácie.
    \item \textbf{PostgreSQL:} Výkonná a robustná databáza s pokročilými funkciami, ideálna pre veľké a zložité aplikácie.
\end{itemize}

(Zdroj: \url{https://www.webglobe.sk/poradna/co-je-databaza})

\subsubsection{Nerelačné databázy (NoSQL)}
Na rozdiel od tabuliek ukladajú dáta flexibilne, najčastejšie ako \textbf{dokumenty}. Tento typ (napr. Firebase Firestore) je kľúčový pre moderné mobilné aplikácie, pretože umožňuje meniť štruktúru dát „za pochodu“ bez nutnosti reštartovania systému \cite{aws_nosql}.
\\ 
(Zdroj: \url{https://aws.amazon.com/nosql/})

\subsubsection{Real-time databázy}
Špeciálny typ databázy, ktorý \textbf{okamžite synchronizuje zmeny} medzi serverom a všetkými pripojenými používateľmi bez nutnosti obnovovania aplikácie. Je to základ pre funkcie ako chaty alebo živé aktualizácie obsahu \cite{firebase_db}.
\\ (Zdroj: \url{https://firebase.google.com/docs/database})

 \fi
\subsection{Správa identít}
\subsubsection{Autentifikácia}
Autentifikácia je proces overenia používateľa teda zistenie, o koho ide. Môžeme overovať podľa kombinácie mena a hesla,
biometrie (napr. otlačku prsta alebo rozpoznania hlasu), alebo iným zariadením 

\subsubsection{Faktory a metódy autentifikácie}
Spôsoby overenia používateľa delíme do troch základných kategórií podľa toho, na čom sú založené:
\begin{itemize}[noitemsep]
    \item \textbf{Znalostný faktor:} Niečo, čo používateľ pozná 
    \item \textbf{Faktor vlastníctva:} Niečo, čo používateľ má 
    \item \textbf{Dedičný faktor:} Niečo, čím používateľ je 
\end{itemize}

Pre dosiahnutie vysokej bezpečnosti sa odporúča viacfaktorové overenie (angl. \textit{Multi-Factor Authentication}, MFA),
ktoré kombinuje aspoň dva z vyššie uvedených faktorov. V mobilných zariadeniach sa čoraz častejšie využíva aj tzv. 
kontinuálne overovanie, ktoré neustále monitoruje identitu používateľa na základe biometrie správania 
 (napr. štýl písania alebo dynamika dotyku) \cite{wiki2025autentifikacia}. 


\iffalse\subsubsection{Digitálna autentifikácia a bezpečnosť}
Proces digitálneho overovania pozostáva z prihlásenia (vytvorenie poverení) a samotnej autentifikácie (predloženie dôkazu identity, napr. tokenu). Pri sieťovej komunikácii je kľúčové chrániť tento proces pred útokmi typu \textit{Man-in-the-Middle} (MITM), kedy sa útočník snaží zachytiť prenášané údaje. 
\\ (Zdroj: \url{https://sk.wikipedia.org/wiki/Autentifikácia})
\fi

\subsubsection{Autorizácia}
Autorizácia je proces overenia oprávnení pre danú úlohu, teda či autentifikovaný používateľ (systém) má oprávnenie vykonať úlohu,
 o ktorú žiada (napr. vstúpiť do systému alebo vykonať platobnú operáciu).
\subsubsection{Mechanizmus trvalého prihlásenia (Token-based Authentication)}
V moderných mobilných aplikáciách je žiaduce, aby používateľ nebol nútený zadávať prihlasovacie
údaje pri každom otvorení aplikácie. Tento proces sa nazýva perzistentná autentifikácia a 
je realizovaný pomocou systému dvoch typov tokenov (bezpečnostných kľúčov):

\begin{itemize}[noitemsep]
    \item \textbf{Access Token (Prístupový kľúč):} Krátkodobý token 
    ktorý mobilná aplikácia pripája ku každej požiadavke na server. Slúži ako dôkaz, že používateľ je momentálne overený.
    \item \textbf{Refresh Token (Obnovovací kľúč):} Dlhodobý token uložený v šifrovanej pamäti zariadenia.
     Jeho jedinou úlohou je vyžiadať si nový \textit{Access Token} v prípade, že pôvodný vypršal.
\end{itemize}

Tento cyklus prebieha na pozadí bez zásahu používateľa (tzv. \textit{Silent Re-authentication}).
Ak je \textit{Refresh Token} stále platný, aplikácia automaticky obnoví reláciu, čím zabezpečí plynulý 
používateľský zážitok pri zachovaní vysokej bezpečnosti, keďže samotné heslo sa sieťou prenáša iba raz – pri prvotnom prihlásení \cite{auth0_refresh_tokens}.


\subsection{Technológia QR kódov a ich štruktúra}
QR kód (\textit{Quick Response Code}) predstavuje dvojrozmerný čiarový kód, ktorý bol navrhnutý pre vysokorýchlostné čítanie a 
veľkú kapacitu dát v porovnaní s tradičnými 1D kódmi. Jeho štruktúra je definovaná niekoľkými kľúčovými prvkami, ktoré umožňujú 
spoľahlivú detekciu aj pri čiastočnom poškodení alebo zlej kvalite obrazu:

\begin{itemize}[noitemsep]
\item \textbf{Vyhľadávacie vzory (Finder Patterns):} Tri identické štvorcové štruktúry v rohoch kódu, ktoré umožňujú skeneru identifikovať polohu a orientáciu kódu v 360-stupňovom rozsahu.
\item \textbf{Korekcia chýb (Error Correction):} Využitie algoritmu \textit{Reed-Solomon}, ktorý umožňuje obnoviť dáta aj v prípade, že je kód znečistený alebo poškodený (v závislosti od úrovne až do 30 % plochy).
\item \textbf{Dátová kapacita:} Na rozdiel od lineárnych kódov dokáže QR kód uchovávať tisíce alfanumerických znakov, čo je kľúčové pre kódovanie komplexných URL adries alebo unikátnych identifikátorov podujatí.
\end{itemize}  \cite{7966807}



  \iffalse
\subsection{Vývoj mobilných aplikácií}


\subsubsection{Mobilná aplikácia}
Mobilná aplikácia je softvérová aplikácia vytvorená špecificky pre smartfóny, tablety alebo iné mobilné zariadenia, ktorá si vyžaduje inštaláciu, čím sa líši od bežného webu. V prvom rade vďaka nim dokážeme rozlíšiť smartfón od bežného telefónu.
\\ (Zdroj: \url{https://www.eliteml.sk/blog/mobilne-aplikacie-v-skratke/})


\subsubsection{Programovacie jazyky}
Programovací jazyk obsahuje inštrukcie vytvorené programátormi, ktoré riadia počítač pri vykonávaní konkrétnych úloh. Tieto inštrukcie, bežne vnímané ako komplexný kód, dodržiavajú odlišnú syntax špecifickú pre každý programovací jazyk a slúžia ako nevyhnutné nástroje pre vývoj softvéru a rôzne výpočtové úlohy.
\\ (Zdroj: \url{https://codelabsacademy.com/sk/blog/what-is-a-programming-language})

\subsubsection{Softvérové frameworky}
Framework je softvérová platforma alebo sada nástrojov, ktorá poskytuje štruktúru a základné funkcionality pre vývoj softvéru. Frameworky uľahčujú vývoj aplikácií tým, že poskytujú preddefinovanú štruktúru, architektúru a sadu nástrojov, ktoré zjednodušujú a urýchľujú proces vývoja.
\\ (Zdroj: \url{https://coderama.com/slovnik/framework})

\subsubsection{Vývojové prostredia a proces spracovania kódu}
Na písanie a správu programovacieho kódu využívajú vývojári špecializované softvérové nástroje, ktoré označujeme ako \textbf{IDE} (\textit{Integrated Development Environment} – Integrované vývojové prostredie). Tieto nástroje siahajú od jednoduchých textových editorov s podporou syntaxe až po komplexné systémy, ktoré obsahujú:

\begin{itemize}
    \item \textbf{Editor kódu:} Miesto na samotný zápis inštrukcií s funkciami ako automatické dopĺňanie kódu.
    \item \textbf{Debugger:} Nástroj na vyhľadávanie a opravu logických chýb v programe.
    \item \textbf{Kompilátor / Interpret:} Špeciálny program, ktorý transformuje (kompiluje) kód napísaný programátorom do binárnej podoby (strojového kódu), ktorej rozumie hardvér mobilného zariadenia alebo procesor.
\end{itemize}

Medzi najpoužívanejšie prostredia pre mobilný vývoj patria napríklad \textbf{Android Studio}, \textbf{Visual Studio Code} (populárne pre frameworky ako React Native alebo Flutter) či \textbf{Xcode} pre platformu iOS. Výber konkrétneho prostredia často závisí od zvoleného programovacieho jazyka a cieľovej platformy.
\\ (Zdroj: \url{https://it-slovnik.cz/pojem/ide-integrovane-vyvojove-prostredi})
%geminy 29.03.2026 8.27 napis mi nieco o prostredi 

  \fi

\subsection{Prístupy k vývoju mobilných aplikácií}
Podľa spoločnosti IBM možno proces vývoja mobilných aplikácií definovať ako tvorbu softvéru navrhnutého 
špecificky pre mobilné koncové zariadenia, pričom kľúčovým rozhodnutím je technológia implementácie. 
V súčasnosti sa vývojári rozhodujú primárne medzi natívnym, multiplatformovým alebo hybridným prístupom \cite{ibm_mobile_app_dev}:

\begin{itemize}[noitemsep]
    \item \textbf{Natívne aplikácie:} Sú vyvíjané priamo pre konkrétnu platformu (Android, iOS)
    \item \textbf{Webové mobilné aplikácie:} Ide o responzívne webové stránky využívajúce technológie HTML5, CSS a JavaScript.
    \item \textbf{Hybridné aplikácie:} Kombinujú predchádzajúce prístupy.
\end{itemize} \cite{olsav2024mobilna}


  \iffalse\subsubsection{Technické obmedzenia a optimalizácia výkonu}
Vývoj pre mobilné platformy vyžaduje špecifický prístup k správe prostriedkov, ktoré sú v porovnaní s desktopovými systémami limitované:

\begin{enumerate}
    \item \textbf{Energetická a systémová efektivita:} Natívne aplikácie vykazujú vyššiu energetickú efektivitu, čo je kľúčové pri aplikáciách využívajúcich lokalizačné služby (GPS) \cite{pullur_medium_2025}.
    \item \textbf{Cloudové spracovanie (\textit{Cloud Offloading}):} Náročné výpočtové operácie, ako sú komplexné odporúčacie algoritmy, je vhodné delegovať na cloudové servery (napr. Firebase).
    \item \textbf{Odozva používateľského rozhrania (UI):} Natívne frameworky dosahujú najlepšiu plynulosť animácií a rýchlosť odozvy na dotyk.
\end{enumerate}

\subsubsection{Distribúcia a ekosystém platforiem}
Platforma Android, s približne 70 \% trhovým podielom, je často preferovaná pre nasadenie MVP kvôli flexibilnejším pravidlám. Naopak, ekosystém iOS vyžaduje striktnejšie dodržiavanie štandardov. Multiplatformové nástroje ako Flutter a React Native umožňujú súčasné nasadenie na oboch platformách s určitým kompromisom v oblasti výkonu.
\\ (Zdroj: \url{https://medium.com/@riyaspullurofficial/native-vs-cross-platform-app-size-and-performance-comparison-550cb7f918b0})


\begin{table}[h!]
\centering
\caption{Súhrnné porovnanie vývojových prístupov (podľa \cite{pullur_medium_2025})}
\label{tab:final_comparison}
\begin{tabular}{|l|c|c|c|c|}
\hline
\textbf{Metrika} & \textbf{Native} & \textbf{Flutter} & \textbf{React Native} & \textbf{KMM} \\ \hline
Veľkosť aplikácie & \textbf{Najmenšia} & Väčšia & Väčšia & Stredná \\ \hline
Zdieľanie kódu & Žiadne & 95 \% + & 90 \% + & Logika \\ \hline
Spotreba batérie & \textbf{Nízka} & Stredná & Vysoká & Nízka \\ \hline
\end{tabular}
\end{table}
  \fi
%\url{https://www.ibm.com/think/topics/mobile-application-development}
%https://medium.com/@riyaspullurofficial/native-vs-cross-platform-app-size-and-performance-comparison-550cb7f918b0 

\subsection{Geolokačné technológie a ich využitie}
Geolokácia je proces identifikácie geografickej polohy zariadenia. V moderných systémoch využívajú štyri hlavné technológie, ktoré sa líšia presnosťou a energetickou náročnosťou:

\begin{itemize}[noitemsep]
    \item \textbf{GPS (Global Positioning System):} Využíva satelitné signály. Je najpresnejší v exteriéri (na metre), ale má vysokú spotrebu batérie a nefunguje vo vnútri budov.
    \item \textbf{Wi-Fi Sniffing:} Určuje polohu skenovaním okolitých Wi-Fi sietí. Spotrebúva menej energie ako GPS.
    \item \textbf{BLE (Bluetooth Low Energy):} Využíva malé vysielače (beacons). Poskytuje extrémnu presnosť v interiéroch (napr. sklady), ale vyžaduje inštaláciu fyzických zariadení v priestore.
    \item \textbf{LPWAN Geolocation:} Využíva nízkoenergetické siete (napr. Sigfox) na trianguláciu signálu. Má najdlhšiu výdrž batérie (roky), ale nižšiu presnosť (stovky metrov).
\end{itemize}\cite{sensolus_geoloc}

V mobilnom vývoji sa najčastejšie využíva hybridná geolokácia, ktorá inteligentne prepína medzi GPS a Wi-Fi podľa dostupnosti
signálu a stavu batérie. To umožňuje aplikáciám efektívne sledovať polohu používateľa bez nadmerného zaťaženia hardvéru.


\subsubsection{Matematické metódy určovania polohy}
Na výpočet presnej polohy zariadenia zo získaných signálov sa v informatike využívajú nasledujúce geometrické a fyzikálne metódy:

\begin{itemize}
    \item \textbf{TOA (Time of Arrival):} Metóda meria absolútny čas, za ktorý signál prejde od vysielača k prijímaču. Na určenie polohy v 2D priestore sú potrebné aspoň tri vysielače (trilaterácia).
    \item \textbf{TDOA (Time Difference of Arrival):} Nameria sa rozdiel v čase príchodu signálov z viacerých vysielačov. Je to praktickejšie ako TOA, pretože nevyžaduje dokonalú synchronizáciu hodín medzi vysielačom a mobilným zariadením.
    \item \textbf{AOA (Angle of Arrival):} Táto metóda neurčuje vzdialenosť, ale uhol, pod ktorým signál prichádza k anténe. Kombináciou uhlov z dvoch staníc je možné určiť presný bod polohy (triangulácia).
    \item \textbf{RSS (Received Signal Strength):} Odhaduje vzdialenosť na základe útlmu sily signálu. Čím je signál slabší, tým je zariadenie pravdepodobne ďalej od vysielača (často využívané pri Wi-Fi a Bluetooth).
\end{itemize} \cite{gentile2012geolocation}


\subsection{Definícia a význam odporúčacích systémov}
V prostredí s takmer neobmedzeným množstvom digitálneho obsahu 
a služieb sa odporúčacie systémy stali kľúčovým nástrojom na personalizovanú filtráciu informácií. 
Ich primárnou úlohou je predpovedať preferencie používateľa a 
zredukovať rozsiahly priestor dostupných položiek na úzky výber
relevantných možností, čím sa zvyšuje celková spokojnosť 
a efektivita interakcie s platformou. \cite{ko2022survey}.

\subsubsection{Obsahové filtrovanie (Content-Based Filtering)}
Tento prístup je založený na analýze špecifických atribútov položiek. Systém identifikuje vlastnosti objektov
(napr. textové informácie, kategórie alebo tagy), o ktoré používateľ v minulosti prejavil záujem, 
a na ich základe generuje odporúčania pre podobné položky \cite{ko2022survey}.
\begin{itemize}[noitemsep]
    \item \textbf{Princíp:} Analýza atribútov a metadát položiek (napr. žáner, kategória) s cieľom nájsť podobnosť s preferenciami používateľa.
    \item \textbf{Využitie:} Efektívne najmä pri službách pracujúcich s textovými dátami, hudbou alebo e-commerce .
    \item \textbf{Limity:} Nízka miera diverzity - používateľovi nie sú ponúkané nové, neobjavené kategórie obsahu, čo vedie k efektu informačnej bubliny.
\end{itemize}\cite{ko2022survey}

\subsubsection{Kolaboratívne filtrovanie (Collaborative Filtering)}
Kolaboratívne filtrovanie predpovedá vkus používateľa na základe analýzy správania a hodnotení 
celej komunity používateľov.

\begin{itemize}[noitemsep]
    \item \textbf{Memory-based:} Hľadá podobnosť priamo medzi používateľmi alebo položkami pomocou algoritmov ako KNN alebo korelačných koeficientov.
    \item \textbf{Model-based:} Využíva pokročilé štatistické metódy a techniky strojového učenia, ako sú zhlukovanie (Clustering), Singular Value Decomposition (SVD) alebo Principal Component Analysis (PCA).
    \item \textbf{Hlavné nedostatky:}  definuje tri kritické obmedzenia:
    \begin{itemize}[noitemsep]
        \item \textit{Cold Start:} Problém s odporúčaním položiek pre nových používateľov bez histórie interakcií.
        \item \textit{Sparsity:} Situácia, keď je v systéme príliš málo hodnotení na to, aby bolo možné spoľahlivo vypočítať podobnosť.
        \item \textit{Gray Sheep:} Používatelia s veľmi unikátnym vkusom, ktorí nevykazujú podobnosť so 
        žiadnou väčšou skupinou používateľov, čo sťažuje generovanie odporúčaní.
    \end{itemize}
\end{itemize}\cite{ko2022survey}

\subsubsection{Hybridné systémy}
Hybridné modely kombinujú viaceré techniky s cieľom eliminovať slabiny jednotlivých prístupov,
najmä problém riedkosti dát (sparsity) a studeného štartu. Podľa Burkeho klasifikácie existuje sedem základných metód hybridizácie, ktoré sú uvedené v tabuľke \ref{tab:hybrid_types}.
\begin{table}[ht]
\centering
\caption{Typy hybridných odporúčacích systémov podľa Burkeho \cite{ko2022survey}}
\label{tab:hybrid_types}
\begin{tabular}{|l|p{9cm}|}
\hline
\textbf{Metóda} & \textbf{Princíp fungovania} \\ \hline
Weighted & Skóre z rôznych techník sa kombinuje pomocou numerických váh. \\ \hline
Switching & Systém prepína medzi technikami v závislosti od kontextu. \\ \hline
Mixed & Odporúčania z viacerých zdrojov sú zobrazené súčasne. \\ \hline
Feature Combination & Atribúty z jedného zdroja slúžia ako vstupné črty pre druhý model. \\ \hline
Feature Augmentation & Výstup jednej techniky sa stáva parametrom pre ďalšie spracovanie. \\ \hline
Cascade & Postupné zjemňovanie poradia (jeden systém vytvorí rebríček, druhý ho upresní). \\ \hline
Meta-level & Celý model vygenerovaný jednou technikou slúži ako vstup pre druhú. \\ \hline
\end{tabular}
\end{table}
Hlavným cieľom väčšiny štúdií zaoberajúcich sa hybridnými modelmi po roku 2010 je integrácia tzv. vedľajších informácií (side information), 
ako sú sociálne väzby medzi používateľmi alebo kontextové dáta, 
ktoré dopĺňajú chýbajúce explicitné hodnotenia. Týmto spôsobom hybridné systémy 
dosahujú vyššiu presnosť a robustnosť v porovnaní so samostatnými modelmi.

\paragraph{Porovnanie systémov}

\begin{table}[H]
\centering
\begin{tabular}{|l|c|c|c|c|}
\hline
\textbf{Kritérium} & \textbf{Collaborative} & \textbf{Content} & \textbf{Hybrid} & \textbf{Knowledge} \\
\hline
Cold Start (používateľ) & Zlé & Zlé & Stredné & Dobré \\
\hline
Cold Start (položka) & Zlé & Dobré & Stredné & Dobré \\
\hline
Škálovateľnosť & Zlé & Dobré & Stredné & Dobré \\
\hline
Diverzita & Dobré & Zlé & Dobré & Stredné \\
\hline
Potreba metadát & Nie & Áno & Čiastočne & Áno \\
\hline
Vysvetliteľnosť & Stredné & Dobré & Stredné & Dobré \\
\hline
\end{tabular}
\caption{Porovnanie typov odporúčacích systémov}
\end{table}

\subsection{Techniky spracovania dát v odporúčacích systémoch}
Pre efektívne fungovanie vyššie spomínaných modelov sa využívajú špecifické techniky dolovania dát (Data Mining) \cite{ko2022survey}:

\begin{enumerate}
    \item \textbf{Text Mining:} Analýza neštruktúrovaných textov pomocou sémantickej analýzy a váhovania pojmov.
    \item \textbf{K-Nearest Neighbor (KNN):} Algoritmus slúžiaci na klasifikáciu položiek alebo používateľov na základe ich blízkosti v metrickom priestore.
    \item \textbf{Clustering (Zhlukovanie):} Automatické rozdeľovanie dát do skupín so spoločnými črtami. Využíva sa najmä pri segmentácii používateľov alebo kategorizácii lokalít.
    \item \textbf{Matrix Factorization:} Pokročilá technika rozkladu matíc, ktorá transformuje interakcie používateľov do priestoru latentných faktorov, čím rieši problém riedkosti dát (Sparsity).
    \item \textbf{Neural Networks:} Modely hlbokého učenia schopné modelovať nelineárne vzťahy a spracovávať komplexné dáta zo sociálnych sietí či senzorov mobilných zariadením.
\end{enumerate}

\iffalse
\subsubsection{Dizajn a používateľská skúsenosť (UX/UI)}

Pri vývoji mobilnej aplikácie nie je kľúčová len funkčná stránka (backend), ale aj spôsob, akým používateľ s aplikáciou interaguje. Návrh používateľského rozhrania sa riadi dvoma základnými disciplínami:

\begin{itemize}
    \item \textbf{UX (User Experience):} Zameriava sa na celkový pocit používateľa pri používaní aplikácie. Cieľom je, aby bola aplikácia intuitívna, logicky usporiadaná a riešila potreby používateľa s minimálnym úsilím.
    \item \textbf{UI (User Interface):} Rieši vizuálnu stránku aplikácie – farby, typografiu, ikony a rozloženie prvkov na obrazovke.
\end{itemize}
%geminy odpoved 29.03.2026 6:36
 \fi


\subsection{Základné princípy moderného návrhu}
Úspešná mobilná aplikácia musí byť postavená na robustnej architektúre, ktorá zaručuje jej 
škálovateľnosť a udržateľnosť v ekosystéme rôznych zariadení.
Pred definovaním konkrétnych vzorov je potrebné určiť pravidlá, ktoré musí moderná architektúra spĺňať pre zachovanie stability v prostredí s obmedzenými systémovými prostriedkami \cite{google_android_architecture}:

\begin{itemize}[noitemsep]
    \item \textbf{Separation of concerns:} Kľúčový princíp, ktorý zakazuje sústredenie kódu v triedach \textit{Activity} alebo \textit{Fragment}. Tieto komponenty slúžia primárne ako hostitelia používateľského rozhrania, sú krátkodobé (\textit{ephemeral}) a operačný systém ich môže kedykoľvek zničiť z dôvodu nedostatku pamäte alebo zmeny konfigurácie.
    \item \textbf{Drive UI from data models:} Používateľské rozhranie by malo byť postavené na dátových modeloch, ktoré sú nezávislé od životného cyklu UI prvkov. Perzistentné modely zabezpečujú, že aplikácia nestratí dáta pri reštarte procesu a zostáva funkčná aj pri nestabilnom pripojení k sieti.
    \item \textbf{Single Source of Truth:} Každý dátový typ má priradeného svojho vlastníka, ktorý ako jediný môže dáta modifikovať. Zmeny sa následne propagujú prostredníctvom imutabilných typov, čo znižuje výskyt chýb a uľahčuje ladenie systému.
    \item \textbf{Unidirectional Data Flow:} Stav aplikácie prúdi smerom nadol (od SSOT k UI elementom) a udalosti/akcie prúdia smerom nahor. Tento vzor zaručuje konzistenciu dát v celom systéme.
\end{itemize}\cite{google_android_architecture}

\subsection{Porovnanie architektonických vzorov}
V histórii vývoja platformy Android sa vyvinulo niekoľko prístupov k organizácii kódu \cite{hostragons2025mvc}:

\begin{enumerate}[noitemsep]
    \item \textbf{MVC (Model-View-Controller):} Tradičný model, kde Controller riadi tok dát. V Androide tento vzor zlyhával, pretože triedy \textit{Activity} preberali úlohu controllera aj view súčasne, čo viedlo k vzniku ťažko udržiavateľného kódu.
    \item \textbf{MVP (Model-View-Presenter):} Oddelil zodpovednosť medzi View a Presenter. Problémom zostala vysoká previazanosť a nutnosť manuálneho ošetrovania životného cyklu, čo spôsobovalo úniky pamäte (\textit{memory leaks}).
    \item \textbf{MVVM (Model-View-ViewModel):} Súčasný štandard. \textit{ViewModel} uchováva stav aplikácie nezávisle od UI komponentov, čím rieši problém s reštartovaním aktivít a stratu dát pri zmene orientácie zariadenia.
\end{enumerate}\cite{hostragons2025mvc}